% ============================================================================
%  EC486 Minor Project — Progress Report
%  SiliconScribe: An AI-Assisted EDA Playground with Autonomous RTL
%  Verification and Self-Correction
%
%  Student : Kushal Pitaliya (23EC105)
%  Institute: CHARUSAT — CSPIT, V.T. Patel Dept. of ECE
%
%  Build   : pdflatex (run twice for page cross-references)
%            pdflatex EC486_Minor_Project_Report.tex
%            pdflatex EC486_Minor_Project_Report.tex
%
%  Style   : Monochrome, minimalist academic (IEEE-Transactions / formal
%            university technical report). Pure grayscale, no colour.
% ============================================================================

\documentclass[11pt,a4paper]{article}

% ---------------------------------------------------------------- layout ----
\usepackage[a4paper,margin=2.2cm,top=2.2cm,bottom=2.2cm,headheight=14pt,headsep=14pt,footskip=18pt]{geometry}
\usepackage[T1]{fontenc}
\usepackage[utf8]{inputenc}
\usepackage{courier}          % monospace for algorithm / code blocks
\renewcommand{\ttdefault}{cmtt}

% --------------------------------------------------------------- graphics ----
\usepackage{amsmath}          % \text{}, math spacing — required for \(\mu\text{m}^2\)
\usepackage{graphicx}
\usepackage{xcolor}           % used ONLY with grayscale values
\usepackage{tikz}
\usetikzlibrary{arrows.meta,positioning,shapes.geometric,calc}

% ----------------------------------------------------------------- tables ----
\usepackage{booktabs}
\usepackage{array}
\usepackage{tabularx}

% ------------------------------------------------------------- typography ----
\usepackage{microtype}
\usepackage{enumitem}
\setlist{nosep,leftmargin=*}
\usepackage{titlesec}
\usepackage{caption}
\captionsetup{font=small,labelfont={bf},justification=justified}
\usepackage{fancyhdr}
\usepackage{lastpage}
\usepackage{hyperref}
\hypersetup{hidelinks,pdftitle={SiliconScribe - EC486 Progress Report},
            pdfauthor={Kushal Pitaliya (23EC105)}}
\usepackage{xurl}

% Paragraph breaking: give TeX a little extra stretch to find a legal line break
% ONLY when a paragraph would otherwise stick out past the margin. This clears
% the small overfull \hbox warnings caused by unbreakable \texttt/\textbf tokens
% without loosening normally-set paragraphs (\sloppy would loosen everything).
\setlength{\emergencystretch}{1.5em}

% ----------------------------------------------------------- colour setup ---
% Monochrome palette only (no saturated hues).
\definecolor{ink}{RGB}{0,0,0}          % pure black
\definecolor{inkdark}{RGB}{31,41,55}   % near-black for body text
\definecolor{rulegrey}{RGB}{107,114,128}
\definecolor{hdrfill}{RGB}{229,231,235}% light grey (headers)
\definecolor{softgrey}{RGB}{243,244,246}
\definecolor{midgrey}{RGB}{156,163,175}

% ------------------------------------------------------------- headings -----
\titleformat{\section}{\normalfont\large\bfseries\color{ink}}
  {\thesection.}{0.6em}{}
\titleformat{\subsection}{\normalfont\normalsize\bfseries\color{ink}}
  {\thesubsection}{0.6em}{}
\titlespacing*{\section}{0pt}{9pt}{4pt}
\titlespacing*{\subsection}{0pt}{6pt}{2pt}

% --------------------------------------------------------- header / footer ---
\pagestyle{fancy}
\fancyhf{}
\renewcommand{\headrulewidth}{0.4pt}
\renewcommand{\footrulewidth}{0.4pt}
\fancyhead[L]{\footnotesize\color{rulegrey}SiliconScribe\ \textbullet\ CHARUSAT \textbullet\ CSPIT ECE}
\fancyhead[R]{\footnotesize\color{rulegrey}EC486: Minor Project}
\fancyfoot[L]{\footnotesize\color{rulegrey}Kushal Pitaliya (23EC105) \textemdash\ Progress Report}
\fancyfoot[R]{\footnotesize\color{rulegrey}Page \thepage\ of \pageref{LastPage}}

% First page style: no running header, keep footer
\fancypagestyle{firstpage}{%
  \fancyhf{}%
  \renewcommand{\headrulewidth}{0pt}%
  \renewcommand{\footrulewidth}{0.4pt}%
  \fancyfoot[L]{\footnotesize\color{rulegrey}Kushal Pitaliya (23EC105) \textemdash\ Progress Report}%
  \fancyfoot[R]{\footnotesize\color{rulegrey}Page \thepage\ of \pageref{LastPage}}%
}

% ----------------------------------------------------------- custom envs ----
\newsavebox{\algoboxsave}
\newsavebox{\codeboxsave}

% LaTeX-style algorithm box (clean top/bottom rules).
\newenvironment{algobox}
  {\par\smallskip\noindent\rule{\linewidth}{0.8pt}\par\vspace{2pt}\small}
  {\par\vspace{2pt}\noindent\rule{\linewidth}{0.8pt}\par\smallskip}

% Code block (for Verilog snippets)
\newenvironment{codebox}
  {\par\medskip\noindent\begin{lrbox}{\codeboxsave}%
   \begin{minipage}{\dimexpr\textwidth-2\fboxsep-2\fboxrule\relax}
   \footnotesize\ttfamily\raggedright}
  {\end{minipage}\end{lrbox}%
   \noindent\fbox{\usebox{\codeboxsave}}\par\medskip}

% Callout: bordered top+bottom with inline lead-in (for the Abstract)
\newcommand{\abstractrule}{\noindent\rule{\textwidth}{0.8pt}\par\vspace{3pt}}

\begin{document}
\thispagestyle{firstpage}

% ============================================================== TITLE =======
\begin{center}
{\color{ink}\rule{\textwidth}{1.0pt}}\\[6pt]
{\large\bfseries CHAROTAR UNIVERSITY OF SCIENCE AND TECHNOLOGY (CHARUSAT)}\\[3pt]
{\normalsize Chandubhai S. Patel Institute of Technology (CSPIT)}\\[2pt]
{\small V.\ T.\ Patel Department of Electronics \& Communication Engineering}\\[4pt]
{\color{ink}\rule{\textwidth}{0.4pt}}\\[10pt]

{\LARGE\bfseries SiliconScribe: An AI-Assisted EDA Playground}\\[4pt]
{\LARGE\bfseries with Autonomous RTL Verification and Self-Correction}\\[12pt]

{\color{ink}\rule{0.55\textwidth}{0.4pt}}\\[10pt]

\noindent\begin{tabular}{@{}>{\bfseries}l@{\hspace{1.5em}}l@{}}
Student Name & Kushal Pitaliya \\
Roll No. & 23EC105 \\
Batch & 2023, B.Tech.\ Electronics \& Communication Engineering \\
Course Code & EC486: Minor Project \\
Academic Year & 2026--2027 (Odd Semester) \\
Submission Type & Progress Report \\
Date of Submission & 15 October 2026 \\
Category & Software-based EDA Tool (Individual Project) \\
\end{tabular}\\[10pt]
{\color{ink}\rule{\textwidth}{1.0pt}}
\end{center}

\vspace{10pt}

% ============================================================== ABSTRACT =====
\abstractrule
\noindent{\textbf{Abstract\ \textemdash}}
In digital VLSI design, writing Register-Transfer Level (RTL) code and developing
self-checking testbenches dominate the front-end of the ASIC development flow, and
verification is widely reported as the single largest share of overall project
effort. While Large Language Models (LLMs) can generate Verilog snippets, they
routinely produce syntax errors, unsupported SystemVerilog constructs, and
functional bugs. More importantly, existing AI code generators operate in an
\emph{open loop}: they cannot execute, simulate, or debug their own code.

To address this, I developed \textbf{SiliconScribe}, a browser-based Electronic
Design Automation (EDA) playground. The system accepts a natural-language
specification, generates synthesizable Verilog-2001 RTL together with a
self-checking testbench, simulates both using Icarus Verilog, and runs an
\textbf{autonomous closed-loop self-correction engine}. When compilation errors or
test failures occur, the simulator log is fed back to the debugger, which emits a
patch and re-simulates until the design passes. To bound the loop, the engine uses
comment- and whitespace-normalised SHA-256 code hashing together with
no-op and oscillation guards. The platform also provides an in-browser SVG digital
waveform viewer, multi-tier linting (Verilator $\rightarrow$ Icarus $\rightarrow$
regex heuristics), Yosys logic-synthesis cell counts, and a 16-file UVM~1.2 bundle
export for QuestaSim. The project is covered by \textbf{42 backend pytest cases and
36 frontend vitest unit tests}. On the development machine 23 backend tests execute
and 19 are skipped because Icarus Verilog is absent from \texttt{PATH}; continuous
integration installs \texttt{iverilog} and \texttt{yosys}, so the full suite runs
there. Known limitations --- notably circular verification, the absence of static
timing analysis, and unmapped synthesis --- are stated explicitly in
Section~\ref{sec:limits}.

\abstractrule
\clearpage

% ========================================================== 1. INTRO ========
\section{Introduction}
\label{sec:intro}

\subsection{Motivation and Problem Statement}
Digital integrated-circuit design demands high precision. Translating a design
requirement into synthesizable RTL, and then proving that RTL correct through a
testbench, is time-consuming and error-prone. Applying general-purpose AI coding
tools directly to hardware design presents three concrete obstacles:

\begin{enumerate}[label=\textbf{\arabic*.}]
\item \textbf{Toolchain incompatibility.} Generic models frequently emit
SystemVerilog constructs (\texttt{logic}, \texttt{always\_ff}, classes,
interfaces) that open-source simulators such as Icarus Verilog cannot compile.
\item \textbf{No execution awareness.} Language models cannot execute code; they
cannot evaluate simulation cycles, clock edges, race conditions, or assertion
matches.
\item \textbf{No closed-loop feedback.} If generated code fails to simulate, an
engineer must manually read the compiler output, locate the bug, and re-prompt the
model --- a slow, manual cycle.
\end{enumerate}

\subsection{Project Scope and Objectives}
This is an individual, software-based minor project for course \textbf{EC486}. The
goal is an autonomous, execution-grounded EDA assistant that bridges natural-language
specification and functional hardware verification. The key objectives are:

\begin{itemize}
\item \textbf{Specification parsing} \textemdash\ convert a prompt into a structured design
schema (\texttt{RTLDesignSpec}).
\item \textbf{Target-constrained generation} \textemdash\ emit Verilog-2001 synthesizable RTL
and a self-checking testbench combining directed corner cases with random stimulus.
\item \textbf{Autonomous debugging loop} \textemdash\ capture diagnostics and assertion
failures, patch the RTL, and re-simulate, with oscillation guards.
\item \textbf{In-browser visualisation} \textemdash\ a native SVG waveform viewer (hex bus
display, zoom, time ruler) plus port-level block diagrams.
\item \textbf{EDA toolchain integration} \textemdash\ multi-tier linting, Yosys synthesis
cell counts, and UVM~1.2 export.
\item \textbf{Zero-budget resilience} \textemdash\ full offline functionality via curated
fallback designs so the system never hard-fails in a demonstration.
\end{itemize}

% ========================================================== 2. WORKING ======
\section{Working: Design, Algorithm, and Methodology}
\label{sec:working}

\subsection{System Architecture}
SiliconScribe uses a modular three-tier architecture (Figure~\ref{fig:arch}).

\begin{figure}[htbp]
\centering
\begin{tikzpicture}[
    font=\footnotesize,
    box/.style={draw=ink,line width=0.6pt,rounded corners=1.5pt,
                fill=softgrey,text width=7.4cm,align=left,
                inner sep=5pt,minimum height=1.0cm},
    arr/.style={-{Latex[length=2.6mm]},line width=0.7pt,color=ink},
    lbl/.style={font=\scriptsize\itshape,color=rulegrey,align=center}
  ]
  \node[box] (fe) {\textbf{FRONTEND}\ \ (React 18 + Vite)\\
    \scriptsize Prompt panel \textbullet\ CodeMirror editor (RTL + TB tabs)
    \textbullet\ Results panel\\
    \scriptsize SVG waveform viewer \textbullet\ Agent chat timeline (SSE)};
  \node[box,below=8mm of fe] (be) {\textbf{BACKEND API}\ \ (FastAPI + Python)\\
    \scriptsize Intent parser $\rightarrow$ RTL generator $\rightarrow$ TB generator\\
    \scriptsize Multi-tier linter \textbullet\ Self-correction loop \textbullet\
    VCD parser\\
    \scriptsize Yosys wrapper \textbullet\ SQLite \textbullet\ UVM generator};
  \node[box,below=8mm of be] (ed) {\textbf{LOCAL EDA TOOLCHAIN}\\
    \scriptsize Icarus Verilog (\texttt{iverilog -g2012 -Wall}, \texttt{vvp})\\
    \scriptsize Yosys (\texttt{synth -top}, \texttt{stat}, \texttt{write\_json})};
  \draw[arr] (fe) -- node[lbl,right=2mm]{HTTP POST / SSE} (be);
  \draw[arr] (be) -- node[lbl,right=2mm]{subprocess} (ed);
\end{tikzpicture}
\caption{Three-tier system architecture of SiliconScribe.}
\label{fig:arch}
\end{figure}

\subsection{End-to-End Pipeline Workflow}
Given a prompt such as \emph{``Design a 4-bit synchronous counter with enable and
async reset''}, the pipeline executes:

\begin{enumerate}[label=\textbf{Step \arabic*:}]
\item \textbf{Intent parsing} \textemdash\ extract module name, parameters, port list
(name/direction/width), behaviour, and constraints into \texttt{RTLDesignSpec}.
\item \textbf{Synthesizable RTL generation} \textemdash\ emit Verilog-2001, enforcing
\texttt{reg}/\texttt{wire} over \texttt{logic} and standard
\texttt{always @(posedge clk)} blocks for Icarus compatibility.
\item \textbf{Pre-simulation linting} \textemdash\ a non-blocking check for syntax errors,
missing semicolons, duplicate modules, and width mismatches.
\item \textbf{Self-checking testbench generation} \textemdash\ instantiate the DUT with named
ports, drive clock/reset, inject \texttt{\$dumpfile}/\texttt{\$dumpvars}, and
compare with \texttt{===} while maintaining pass/fail counters.
\item \textbf{Simulation} \textemdash\ compile with \texttt{iverilog -o testbench.vvp
-g2012 -Wall} and run via \texttt{vvp} in an isolated workspace under a timeout.
\item \textbf{Autonomous debugging} \textemdash\ if the result is FAIL or ERROR, invoke the
self-correction loop (Section~\ref{sec:loop}).
\item \textbf{Post-processing} \textemdash\ evaluate pass rate, parse \texttt{design.vcd}
into waveform data, run Yosys for cell counts, persist to SQLite, and stream results
to the frontend.
\end{enumerate}

\subsection{The Autonomous Closed-Loop Debugger}
\label{sec:loop}
The core contribution is the self-correction loop (Algorithm 1).
Hash normalisation strips single-line comments (\texttt{//}), block comments
(\texttt{/* */}), and all whitespace before SHA-256 is computed. A candidate patch
whose hash is already in \texttt{SeenHashes} terminates the loop immediately,
which both prevents oscillation (the model flipping between two buggy variants) and
conserves tokens. A \emph{no-op} guard separately catches the case where the model
returns the input unchanged.

\begin{algobox}
{\normalsize\bfseries\rmfamily Algorithm 1\ \ \ Autonomous Hardware Self-Correction}\\[2pt]
{\rmfamily\rule{\linewidth}{0.4pt}}\\[2pt]
{\rmfamily\textbf{Input:} RTL$_0$, TB$_0$, MaxIterations\\
\textbf{Output:} RTL, TB, Status}\\[3pt]
{\rmfamily
\begin{tabularx}{\linewidth}{@{}r@{\hspace{0.75em}}>{\raggedright\arraybackslash}X@{}}
1.  & RTL $\gets$ RTL$_0$;\quad TB $\gets$ TB$_0$ \\
2.  & SeenHashes $\gets$ \{ NormHash(RTL$_0$) \} \\
3.  & Status, Log, Pass, Fail $\gets$ Simulate(RTL, TB) \\
4.  & Iter $\gets$ 0 \\
5.  & \textbf{while} Status $\in$ \{FAIL, ERROR\} \textbf{and} Iter $<$ MaxIterations \textbf{do} \\
6.  & \hspace*{1.2em} Iter $\gets$ Iter $+$ 1 \\
7.  & \hspace*{1.2em} Patch $\gets$ LLM\_FixDesign(RTL, TB, Log) \\
8.  & \hspace*{1.2em} NewRTL $\gets$ Patch.rtl\_code \\
9.  & \hspace*{1.2em} NewHash $\gets$ NormHash(NewRTL) \\
10. & \hspace*{1.2em} \textbf{if} NewHash $=$ NormHash(RTL) \textbf{then} \ \textbf{return} ``no-op'' \\
11. & \hspace*{1.2em} \textbf{if} NewHash $\in$ SeenHashes \textbf{then} \ \textbf{return} ``oscillation'' \\
12. & \hspace*{1.2em} SeenHashes $\gets$ SeenHashes $\cup$ \{NewHash\} \\
13. & \hspace*{1.2em} RTL $\gets$ NewRTL \\
14. & \hspace*{1.2em} Status, Log, Pass, Fail $\gets$ Simulate(RTL, TB) \\
15. & \textbf{end while} \\
16. & \textbf{return} RTL, TB, Status \\
\end{tabularx}}
\end{algobox}

\subsection{Digital Waveform Parsing and SVG Rendering}
Rather than requiring an external viewer such as GTKWave, SiliconScribe embeds a
VCD parser (\texttt{backend/vcd\_parser.py}). It reads time steps
(\texttt{\#0}, \texttt{\#10}, \ldots) and value transitions, converts multi-bit
vectors into hexadecimal (\texttt{0x0}, \texttt{0xA}) while preserving unknown
(\texttt{x}) and high-impedance (\texttt{z}) states, and enforces safety limits
($\le 40$ signals, $\le 2000$ changes per signal) to prevent browser memory
exhaustion. The frontend (\texttt{WaveformViewer.tsx}) renders responsive SVG with
zoom, pan, and a dynamic time ruler.

\subsection{UVM~1.2 Export Engine}
SiliconScribe generates a 16-file UVM testbench suite
(\texttt{backend/uvm\_templates.py}): a DUT interface with clocking blocks, sequence
items, sequences, sequencers, drivers, monitors, a scoreboard with an analysis
imp, a functional-coverage subscriber, a testbench top, a package file, and a
QuestaSim \texttt{Makefile} and \texttt{filelist.f}. The bundle is packaged
in-memory and delivered as a downloadable \texttt{.zip}. It is an
\textbf{export-only} artefact --- see Limitation~\ref{lim:uvm}.

% ========================================================== 3. RESULTS ======
\section{Results and Discussion}
\label{sec:results}

\subsection{Work Completed (Milestone: 15/10/2026)}
\begin{itemize}
\item Full end-to-end pipeline from natural-language prompt to simulation output.
\item Autonomous self-correction engine with verified hash-based cycle prevention.
\item Multi-tier linter and Yosys logic-synthesis wrapper.
\item Interactive web frontend: CodeMirror editor, SVG waveform viewer, project
history drawer.
\item SQLite persistence and the 16-file UVM bundle generator.
\item 100\% offline fallback mode with 5 pre-validated hardware design modules.
\end{itemize}

\subsection{Case Study: Self-Correcting 4-Bit Up-Counter}
A deliberately buggy prompt was submitted:
\begin{quote}
\itshape ``Design a 4-bit synchronous up-counter with async reset and enable, but
make it buggy.''
\end{quote}

\noindent\textbf{Mode disclosure (important).} SiliconScribe runs in two provider
modes, and the loop behaves differently in each. This case study ran in
\textbf{offline demo mode} (\texttt{OFFLINE\_MODE=1}) --- the mode used by all
automated regression tests and by the deployed Docker image when no API key is
present. In offline mode the system serves \emph{curated} designs: the buggy RTL
(\texttt{count <= count + 2'd2}) is a known defect seeded deliberately, and the
``fix'' is a \textbf{scripted swap to a stored known-correct implementation},
reported honestly in the UI timeline and in \texttt{iteration\_history} as
\texttt{fix\_type = "offline\_scripted"}. Its purpose is to validate the
\textbf{loop machinery} --- defect detection, log parsing, re-simulation, normalised
hash cycle prevention, and the PASS/FAIL decision path --- deterministically and at
zero API cost. In \textbf{live LLM mode} (a provider key configured), the same loop
instead sends RTL, testbench, and log to the model, which performs genuine
log-driven diagnosis. The study below is therefore a \emph{fault-injection
regression test of the loop}, not a claim that a language model diagnosed the bug.

\begin{itemize}
\item \textbf{Iteration 0 (initial attempt):} the served RTL carried the seeded
defect \texttt{count <= count + 2'd2;}. The testbench ran 20 directed vectors:
\textbf{0 Passed, 20 Failed} (status FAIL). Transcript:
\texttt{FAIL: expected=1 got=2}, \texttt{FAIL: expected=2 got=4}, \ldots
\item \textbf{Iteration 1 (automated patch):} the loop captured the failure log and
invoked the debugger; the known-correct implementation was substituted
(\texttt{count <= count + 1'b1;}), tagged \texttt{offline\_scripted}. Re-simulation
returned \textbf{20 Passed, 0 Failed} (status PASS). The UI displayed the full
resolution timeline including the scripted-fix disclosure.
\end{itemize}

\subsection{Offline Benchmark Designs}
Five designs were implemented and benchmarked to guarantee reliability without
internet access (Table~\ref{tab:bench}).

\begin{table}[htbp]
\centering
\caption{Offline benchmark designs and simulation results.}
\label{tab:bench}
\footnotesize
\setlength{\tabcolsep}{4pt}
\begin{tabular}{@{}l l c r r c@{}}
\toprule
\textbf{Module} & \textbf{Functionality} & \textbf{Width} & \textbf{Time}
  & \textbf{Vectors} & \textbf{Result} \\
\midrule
\texttt{counter}   & Up-counter, enable, async reset & 4b & 100 ns
  & 20 (dir.) & 100\% PASS \\
\texttt{alu\_4bit} & ALU: ADD SUB AND OR XOR          & 4b & Comb
  & 200 (dir+rand) & 100\% PASS \\
\texttt{alu\_8bit} & 8-bit ALU, carry/overflow        & 8b & Comb
  & 200 (dir+rand) & 100\% PASS \\
\texttt{full\_adder} & Gate-level full adder          & 1b & Comb
  & 8 (exhaust.) & 100\% PASS \\
\texttt{mux4to1}   & 4-to-1 mux, parameterised        & 8b & Comb
  & 100 (dir+rand) & 100\% PASS \\
\bottomrule
\end{tabular}
\end{table}

Vector counts are taken from the testbench loop bounds in
\texttt{backend/\allowbreak offline\_designs.py}: \texttt{repeat (20)} for the counter;
\texttt{for i = 0..199} for both ALUs; \texttt{for i = 0..7} (exhaustive) for the
full adder; and \texttt{for i = 0..99} for the mux. The counter testbench is fully
directed (no \texttt{\$random}); the ALU and mux testbenches mix directed coverage
with randomized stimulus.

\subsection{Logic Synthesis Metrics (Yosys)}
Yosys produced genuine standard-cell counts. \textbf{An important caveat applies:}
Yosys is invoked \emph{without} a liberty file, so RTL is mapped onto Yosys's
generic gate library (\texttt{\$\_AND\_}, \texttt{\$\_NOT\_}, \texttt{\$\_DFF\_},
\ldots). The flow therefore performs \textbf{no technology mapping} and yields
\textbf{no physical area figure}. The \texttt{area\_estimate} field returned by the
API is a \textbf{dimensionless relative proxy} (\texttt{cell\_count}~$\times$~5)
used only to rank designs against one another --- it is explicitly \emph{not} an
area in \(\mu\mathrm{m}^2\). The API flags this via
\texttt{liberty\_mapped: false}, \texttt{area\_units: "relative\_proxy"}, and
\texttt{tech\_node: null}.

\begin{table}[htbp]
\centering
\caption{Yosys synthesis metrics (unmapped, generic gate library).}
\label{tab:synth}
\small
\begin{tabular}{@{}lrrl@{}}
\toprule
\textbf{Design} & \textbf{Generic cells} & \textbf{Relative area proxy}
  & \textbf{Physical area} \\
\midrule
\texttt{full\_adder} & 2  & 10  & \textemdash\ (no liberty) \\
\texttt{counter}     & 12 & 60  & \textemdash\ (no liberty) \\
\texttt{alu\_4bit}   & 22 & 110 & \textemdash\ (no liberty) \\
\bottomrule
\end{tabular}
\end{table}

Obtaining genuine PPA figures requires a PDK liberty file (e.g.\ SkyWater 130\,nm
or Nangate 45\,nm) plus technology mapping and static timing analysis; this is
planned for the final submission.

\subsection{Automated Test Suites}
\begin{itemize}
\item \textbf{Backend (pytest):} \textbf{42 test cases} covering Pydantic models,
linter heuristics, VCD parsing, rate limiting, and the orchestrator loop. On the
development machine \textbf{23 pass and 19 skip} because Icarus Verilog is not on
\texttt{PATH} (\texttt{skipif(shutil.which("iverilog") is None)}). CI installs
\texttt{iverilog} and \texttt{yosys}, so all 42 execute there.
\item \textbf{Frontend (vitest):} \textbf{36 unit tests} covering
\texttt{PromptPanel}, \texttt{ResultsPanel}, and \texttt{WaveformViewer}.
\item \textbf{TypeScript build:} zero errors (\texttt{tsc -b \&\& vite build}).
\item \textbf{Continuous integration:} \texttt{ruff} lint, the full pytest suite
with tools installed, an end-to-end \texttt{POST /api/\allowbreak design/run} assertion
requiring \texttt{status == "PASS"}, a 16-file UVM generation check, frontend
type-check/test/build, and a Docker image build.
\end{itemize}

\noindent\textit{Honest note on test strength.} 19 of the 23 locally executed
backend tests exercise pure-Python paths (models, linter fallbacks, VCD parsing,
rate limiting, sanitisation). The tests that prove the self-correction loop reaches
a passing state require a real simulator, which is why they skip locally and run in
CI.

% ====================================================== 4. LIMITATIONS ======
\section{Limitations}
\label{sec:limits}

\begin{enumerate}[label=\textbf{L\arabic*.},ref=L\arabic*]
\item \textbf{Circular verification (most important limitation).\ }
\label{lim:circular}The same language model generates both the RTL \emph{and} the
self-checking testbench, and the debugger prompt instructs it to ``treat the
testbench's intent as the specification.'' The loop therefore demonstrates
\textbf{self-consistency, not ground-truth correctness}: if the model misreads the
specification (e.g.\ swapping active-high and active-low reset), it can emit a DUT
and a testbench that are \emph{consistently wrong together} and still report 100\%
PASS. Closing this requires an \textbf{independent golden reference model} (e.g.\
a Python behavioural model the testbench compares against), planned for the final
submission.
\item \textbf{No static timing analysis or PPA sign-off.}
\label{lim:sta}The \texttt{target\_frequency\_mhz} field is appended to the
RTL-generation prompt as a textual hint only. There is no SDC, no
\texttt{create\_clock}, no \texttt{abc -D}, and no setup/hold analysis in the flow.
The platform reports no timing slack and no real power/area/performance figures.
\item \textbf{Synthesis is unmapped (no technology library).}
\label{lim:mapping}Yosys runs without a liberty file, so RTL maps to the generic
gate library. Cell counts are genuine, but \texttt{area\_estimate} is a
dimensionless relative proxy, not physical area (Section~\ref{sec:results}).
\item \textbf{Testbench-reported, not instrumented, coverage.}
\label{lim:cov}The coverage figure is a regex parse of \texttt{\$display} output
emitted by the testbench itself. It is not instrumented toggle/line/branch
coverage, which would require simulator instrumentation such as
\texttt{verilator {-}{-}coverage}.
\pagebreak
\item \textbf{UVM bundle is export-only and never compiled.}
\label{lim:uvm}The 16-file UVM~1.2 bundle is generated as SystemVerilog
(\texttt{logic}, \texttt{always\_ff}, classes) for a commercial simulator. Icarus
Verilog cannot compile UVM, and no commercial licence was available, so the bundle
has \textbf{not been compiled or simulated}. It is a structured template export,
not a verified working testbench; CI checks structural integrity (file count,
generation success), not compilation.
\item \textbf{Schematic display depth.} The browser shows a port-level block
diagram. Gate-level logic diagrams require the planned \texttt{netlistsvg}
integration over Yosys JSON.
\item \textbf{Design hierarchy scope.} The implementation targets single-module or
small blocks ($\le 1000$ RTL lines). Multi-clock SoC architectures exceed
single-prompt LLM context windows.
\end{enumerate}

% ======================================================= 5. CONCLUSION ======
\section{Conclusion and Next Steps}

\subsection{Conclusion}
For this progress milestone, \textbf{SiliconScribe} has been designed,
implemented, and partially verified. Combining natural-language processing,
synthesizable Verilog-2001 generation, Icarus Verilog simulation, and an autonomous
self-correction loop with oscillation guards, the project demonstrates that an LLM
can act as an execution-grounded assistant for digital design and verification ---
while its current boundaries (circular verification, absent STA, unmapped
synthesis) are documented rather than hidden.

\subsection{Plan for Final Submission (Deadline: 06/11/2026)}
In priority order:
\begin{enumerate}[label=\textbf{\arabic*.}]
\item \textbf{Independent golden reference model (highest priority).} Add a
language-independent behavioural model (Python) for each benchmark so the testbench
compares the DUT against the model rather than against an expectation derived from
the same RTL, breaking circular verification.
\item \textbf{Real PPA via liberty-based technology mapping.} Run Yosys with an
open PDK liberty file (\texttt{read\_liberty}, \texttt{abc}, \texttt{stat -tech}),
replacing the area proxy with genuine cell area and estimated delay.
\item \textbf{Gate-level netlist visualisation} via \texttt{netlistsvg} over Yosys
JSON.
\item \textbf{Instrumented coverage} via \texttt{verilator {-}{-}coverage}.
\item \textbf{Poster and presentation} for the final EC486 evaluation.
\end{enumerate}

% ======================================================== 6. REFERENCES =====
\section*{References}
\begin{enumerate}[label={[\arabic*]},leftmargin=2.2em,itemsep=0pt,parsep=0pt]
\item IEEE Standard for SystemVerilog --- Unified Hardware Design, Specification,
and Verification Language, \textbf{IEEE Std 1800-2017}, 2017.
\item Accellera Systems Initiative, \textit{Universal Verification Methodology
(UVM) 1.2 User's Guide}, Accellera, 2015.
\item S.~Thakur et al., ``VeriGen: A Large Language Model for Verilog Code
Generation,'' \textit{ACM TODAES}, vol.~29, no.~3, 2024.
\item M.~Liu et al., ``ChipNeMo: Domain-Adapted LLMs for Chip Design,'' in
\textit{Proc.~61st ACM/IEEE Design Automation Conference (DAC)}, 2024.
\item S.~Williams, ``Icarus Verilog: An Open-Source Verilog Compiler and
Simulator,'' \url{http://iverilog.icarus.com/}.
\item C.~Wolf, ``Yosys Open Synthesis Suite,'' in \textit{Proc.~21st Int.\ Workshop
on Logic \& Synthesis (IWLS)}, 2012.
\item W.~Snyder, ``Verilator: Fast Open-Source Verilog Simulator,''
Veripool.org, 2024.
\end{enumerate}

\vspace{4pt}
\noindent{\color{ink}\rule{\textwidth}{0.8pt}}
\begin{center}
\small\itshape End of Progress Report
\end{center}
\noindent{\color{ink}\rule{\textwidth}{0.8pt}}

\end{document}
